\editionPage[\pdfbookmark[0]{Corrigés des sujets blancs}{bm-corriges-sujets}]{sujet-A-cor-auto}{Sujet A · Corrigé}{Automatismes : réponses et méthode}{\hyperref[sujet-A]{Énoncé p.~\pageref*{sujet-A}}}
\begin{blocEdition}{sujet-A-cor-auto}
\begin{questions}\item $60\times1{,}15=\res{69\text{ euros}}$.
\item On divise par le coefficient multiplicateur : $75/1{,}25=\res{60\text{ euros}}$.
\item Le coefficient réciproque est $1/0{,}8=1{,}25$ : \res{25\,\%}.
\item Par différence de deux carrés : $4-3=\res{1}$.
\item Un produit est nul si l’un de ses facteurs est nul : $\res{S=\{-3\,;\,1/2\}}$.
\item On divise par $-3$ en changeant le sens : \res{x<2}, soit $\mathopen{]}-\infty;2[$.
\item Le coefficient de $x^2$ vaut 1 : le produit est strictement négatif entre ses racines $-4$ et 1, donc $\res{x\in\mathopen{]}-4;1[}$.
\item $m=\dfrac{-5-3}{3-(-1)}=\dfrac{-8}{4}=\res{-2}$.
\item $P(A\cap B)=P(A)\times P_A(B)=0{,}4\times0{,}25=\res{0{,}1}$.
\item La série ordonnée a 5 valeurs ; la troisième est la médiane : \res{2}.
\item La courbe est sur ou sous la droite d’équation $y=1$ pour les abscisses de $-1$ à 3 : $\res{S=[-1;3]}$.
\item L’ordonnée à l’origine vaut 3 ; de $(0;3)$ à $(2;2)$, $y$ diminue de 1 quand $x$ augmente de 2 : $\res{y=-\tfrac12x+3}$.\end{questions}
\label{figedition-sujet-A-cor-auto-1}\figureGeo[corrige]{fr-sa-auto}
\end{blocEdition}

\editionPage{sujet-A-cor-1}{Sujet A · Corrigé 1}{Une courbe et sa tangente}{\hyperref[sujet-A-ex-1]{Énoncé p.~\pageref*{sujet-A-ex-1}} · Raisonnement et rédaction}
\begin{blocEdition}{sujet-A-cor-1}

\editionEtape{1. Dérivée, signes et variations}
On obtient $f'(x)=3x^2-12x+9=3(x-1)(x-3)$.
Sur $[0;4]$, les signes sont successivement $+$, $-$, $+$, avec des zéros en 1 et 3.
\begin{center}\begin{tableau}\otTab[2.1]{$x$/1,$f'(x)$/1,$f$/2}{0,1,3,4}
\tkzTabLine{,+,z,-,z,+,}\tkzTabVar{-/2,+/6,-/2,+/6}\end{tableau}\end{center}
Le minimum est \res{2}, atteint en 0 et 3 ; le maximum est \res{6}, atteint en 1 et 4.
Les bornes du segment sont des candidats aux extremums globaux et doivent être comparées aux points critiques.
\editionEtape{2. Tangente et position relative}
On a $f(2)=4$ et $f'(2)=-3$. Ainsi $T:y=-3(x-2)+4=-3x+10$.
Le développement de $(x-2)^3$ donne $x^3-6x^2+12x-8$, égal à $f(x)-(-3x+10)$.
Ce cube est négatif pour $x<2$, nul en 2, positif pour $x>2$.
La courbe est donc sous la tangente sur $[0;2[$ et au-dessus sur $]2;4]$.
Elle la traverse au point de contact : une tangente ne reste pas forcément du même côté de la courbe.

\label{figedition-sujet-A-cor-1-1}\figureGeo[corrige]{fr-sa-fonction}
\end{blocEdition}

\editionPage{sujet-A-cor-2}{Sujet A · Corrigé 2}{Ateliers, défauts et gain}{\hyperref[sujet-A-ex-2]{Énoncé p.~\pageref*{sujet-A-ex-2}} · Raisonnement et rédaction}
\begin{blocEdition}{sujet-A-cor-2}

\editionEtape{1. Reconstituer les probabilités}
Les poids manquants sont $19/20$ et $9/10$. Par probabilités totales,
\[P(D)=\frac35\frac1{20}+\frac25\frac1{10}=\frac3{100}+\frac4{100}=\res{\frac7{100}}.\]
Les articles défectueux provenant de B représentent $4/100$ de tous les articles.
En se limitant aux défauts, on obtient
$P_D(B)=(4/100)/(7/100)=\res{4/7}$.
\editionEtape{2. Passer de la loi aux deux moments}
Le gain $G$ prend $-8$ avec probabilité $7/100$ et 12 avec probabilité $93/100$.
Ainsi $E(G)=(-8)(7/100)+12(93/100)=\res{53/5=10{,}6}$.
Le deuxième moment vaut $64(7/100)+144(93/100)=692/5$.
Il vient
\[V(G)=\frac{692}{5}-\left(\frac{53}{5}\right)^2=\res{\frac{651}{25}=26{,}04}.\]
\editionEtape{3. Utiliser l'évènement contraire}
Deux articles sans défaut ont pour probabilité $(93/100)^2$ grâce à l'indépendance.
La probabilité demandée est donc $1-(93/100)^2=\res{1351/10000=0{,}1351}$.
L'évènement « au moins un » comprend un ou deux défauts ; compter seulement un défaut serait incomplet.

\label{figedition-sujet-A-cor-2-1}\figureGeo[corrige]{fr-sa-arbre}
\end{blocEdition}

\editionPage{sujet-B-cor-auto}{Sujet B · Corrigé}{Automatismes : réponses et méthode}{\hyperref[sujet-B]{Énoncé p.~\pageref*{sujet-B}}}
\begin{blocEdition}{sujet-B-cor-auto}
\begin{questions}\item Au dénominateur 12 : $\dfrac9{12}+\dfrac2{12}=\res{\dfrac{11}{12}}$.
\item $(2^3)^2=2^6$, donc le quotient vaut $2^{6-4}=\res{4}$.
\item Différence de deux carrés : $(3x)^2-1^2=\res{(3x-1)(3x+1)}$.
\item $\dfrac{92-80}{80}=\dfrac{12}{80}=0{,}15$ : hausse de \res{15\,\%}.
\item $(1+t)^2=1{,}21=1{,}1^2$ et $1+t>0$, donc \res{t=10\,\%}.
\item On divise par $\pi r^2\neq0$ : $\res{h=\dfrac{V}{\pi r^2}}$.
\item L’ordonnée à l’origine est l’ordonnée de $A$, soit 3, et la pente vaut $-2$ : $\res{y=-2x+3}$.
\item $\dfrac{8+2\times12+15}{1+2+1}=\dfrac{47}4=\res{11{,}75}$.
\item Parmi les 120 filles, 30 sont internes : $P_F(I)=\dfrac{30}{120}=\res{\dfrac14}$.
\item $P(F\cap I)=\dfrac{30}{200}=\res{\dfrac3{20}}$ ; parmi les 60 internes, 30 sont des filles : $P_I(F)=\res{\dfrac12}$.
\item $g$ s’annule en $-1$ et en 3 ; \res{g(x)>0} sur $\mathopen{]}-1;3[$ et \res{g(x)<0} sur $[-2;-1[$ et sur $]3;4]$.
\item La médiane vaut \res{12} ; l’écart interquartile vaut $Q_3-Q_1=15-9=\res{6}$.\end{questions}
\label{figedition-sujet-B-cor-auto-1}\figureGeo[corrige]{fr-sb-auto}
\end{blocEdition}

\editionPage{sujet-B-cor-1}{Sujet B · Corrigé 1}{Une évolution avec apport}{\hyperref[sujet-B-ex-1]{Énoncé p.~\pageref*{sujet-B-ex-1}} · Raisonnement et rédaction}
\begin{blocEdition}{sujet-B-cor-1}

\editionEtape{1. Déplacer la valeur d'équilibre}
On calcule $v_{n+1}=u_{n+1}-40=(3/4)(u_n-40)=(3/4)v_n$ et $v_0=40$.
Ainsi $v_n=40(3/4)^n$ et \res{u_n=40+40(3/4)^n}.
La puissance est positive et strictement décroissante ; on a $u_n>40$ et la suite décroît strictement.
\editionEtape{2. Prouver le premier seuil}
L'inégalité $u_n\leq50$ équivaut à $(3/4)^n\leq1/4$.
Au rang 4, $81/256>64/256=1/4$ ; au rang 5, $243/1024<256/1024=1/4$.
La décroissance établie précédemment prouve donc que le premier rang est \res{5}.
\begin{python}
def premier_rang():
    n, u = 0, 80.0
    while u > 50:
        u = 0.75 * u + 10
        n = n + 1
    return n
\end{python}
Le test \code{u > 50} continue tant que le seuil n'est pas atteint.
À la sortie, $u=u_n\leq50$ ; le compteur correspond au nombre d'actualisations effectuées.

\end{blocEdition}

\editionPage{sujet-B-cor-2}{Sujet B · Corrigé 2}{Cercle, tangente et projection}{\hyperref[sujet-B-ex-2]{Énoncé p.~\pageref*{sujet-B-ex-2}} · Raisonnement et rédaction}
\begin{blocEdition}{sujet-B-cor-2}

\editionEtape{1. Construire le centre du cercle}
La médiatrice de $[AB]$ est $x=2$.
Pour un point $(x;y)$, l'égalité des distances à $A$ et à $C$ s'écrit
$x^2+y^2=(x-1)^2+(y-3)^2$, soit $x+3y=5$.
L'intersection des deux médiatrices est \res{\Omega(2;1)}.
Le rayon au carré est $\Omega A^2=5$ : le cercle est
\[\res{(x-2)^2+(y-1)^2=5}.\]
Pour $C(1;3)$, le membre de gauche vaut $1+4=5$ : l'appartenance est vérifiée.
\editionEtape{2. Utiliser la direction normale du rayon}
Le rayon $\vect{\Omega C}=(-1;2)$ est normal à la tangente en $C$.
Son équation est $-(x-1)+2(y-3)=0$, donc \res{-x+2y-5=0}.
\editionEtape{3. Projeter et contrôler}
Le projeté s'écrit $H=M-\lambda(-1;2)=(5+\lambda;4-2\lambda)$.
L'équation de la tangente impose $-(5+\lambda)+2(4-2\lambda)-5=0$,
soit $-2-5\lambda=0$. Ainsi \res{H(23/5;24/5)}.
On vérifie $-23/5+48/5-5=0$. De plus $\vect{MH}=(-2/5;4/5)$ est parallèle au normal.
Les deux conditions du projeté sont donc satisfaites et
\[MH=\sqrt{4/25+16/25}=\res{\frac{2\sqrt5}{5}}.\]

\label{figedition-sujet-B-cor-2-1}\figureGeo[corrige]{fr-sb-geometrie}
\end{blocEdition}

